\begin{abstract}
Large language models are increasingly used as annotators in computational social science, yet annotation errors propagate into downstream causal estimates in ways that depend on both the confusion matrix structure and the causal DAG topology. We derive bias expressions for seven canonical topologies under $K$-class non-differential misclassification and show that $K {\geq} 3$ introduces a structural possibility of sign reversal absent in the binary case. We prove sign preservation for confounding and collider, and verify it numerically for mediation, M-bias, and front-door; exposure and IV topologies can reverse the treatment effect sign. Under realistic LLM accuracy, sign reversal occurs in 2.1\% of exposure configurations, in secondary coefficients with small effect sizes; across all 23 empirical LLM configurations (mean diagonal 0.45--0.81), sign reversal is confined exclusively to exposure, though IV reaches 52.6\% under uniform Dirichlet priors. We introduce Annotation Sensitivity Values (ASVs), with proved ($K{=}2$) and numerically verified ($K{\geq}3$) monotonicity, yielding a vulnerability ranking from IV (maximally fragile) to M-bias (immune), and show that matrix-based corrections are provably ineffective under linear regression for four topologies due to Frisch--Waugh--Lovell invariance. Topology-aware analysis is necessary for valid causal inference with LLM-generated annotations.
\end{abstract}
